\newif\ifglobal

%%%%%%%%%%%%%%%%%%%%%%%
%%% Compile to view %%%
%%%%%%%%%%%%%%%%%%%%%%%

%%%%%%%%%%%%%%%%%%%%%%%%%%%%%%%%%%%%%%%%%%%%%%%%%%%
%%% Readme:                                     %%%
%%% https://www.overleaf.com/read/bwdjvfjksvjy/ %%%
%%%%%%%%%%%%%%%%%%%%%%%%%%%%%%%%%%%%%%%%%%%%%%%%%%%

% >>> M?: marks a module setting number or important notice

% >>> M4: Set {./relative-path} to external files for this module <<<
\def \modulefiles {./04-DMA}
% To add an external file, use:
% (Replace '---' with actual filename)
% '\input{\modulefiles/tables/---.tex}' for tables
% '\includegraphics{\modulefiles/figures/---}' for figures
% '\subfile{\modulefiles/---}' for register files


% >>> M1: This module is in global project? <<<
% 'YES' - uncomment; 'NO' - comment out
\globaltrue

\ifglobal
    \documentclass[main\_\_CN.tex]{subfiles}

\else
    % Font "FandolSong-Regular" used by ctex does not have GSUB table.
    % It causes warnings that do not affect compilation and can be ignored.
    \PassOptionsToPackage{quiet}{fontspec}

    \documentclass[openany, 10pt]{book}

    % >>> M2: In ./chip-spec-settings/preamble-trm-module, also find
    %         settings for visibility of todo notes, watermark <<<
    \usepackage{preamble-trm-module}


    % Enable Chinese typesetting
    \usepackage{ctex}

    % >>> M3: Temporary labels in Chinese
    %         you can add your own labels there <<<
    \myexternaldocument{temp-labels-cn}

\fi %\ifglobal


\setcounter{secnumdepth}{4}

% Define formatting of chapters in text
\titleformat{\chapter}
  [display]
  {\LARGE\bfseries}
  {\color{AllportsColor} 第\, \thechapter\, 章}
  {1em}
  {\color{AllportsColor}}
\titlespacing*{\chapter}{0pt}{0pt}{20pt}

% Define headers
\usepackage{fancyhdr}
\usepackage{lipsum}
\renewcommand{\chaptermark}[1]{\markboth{第\, \thechapter\, 章 \  \ #1}{}}
\fancypagestyle{plain}{%
  \fancyhead{}%
  %\fancyhead[R]{\Navigation}
  \fancyhead[L]{%
    \nouppercase{%
      \ %
      \normalsize\leftmark
    }%
  }%
}
\pagestyle{plain}

% >>> M5: If this module needs any updates in
% './chip-spec-settings/preamble-trm-module.sty'
% (include additional package, etc.), contact your module PM
% or create the file './chip-spec-settings/changelog.tex' make the
% necessary updates yourself and document them in this file <<<

% Make text left-aligned and allow latex to break pages naturally, comment out for full justification
\raggedright
\raggedbottom


\begin{document}


% Sets language into which pre-defined bits of text are translated
\selectlanguage{chinese}

% Do not delete this block
\ifglobal
% Add the following front matter if
% this module is outside of global project
\else
    \InputIfFileExists{readme}{}{}
    \listoftodos
    {\let\clearpage\relax \listofdonetodos}
    \tableofcontents
    %\listoftables
    %\listoffigures
\fi


%%%%%%%%%%%%%%%%%%%%%%%%%
%%% TRM Module Begins %%%
%%%%%%%%%%%%%%%%%%%%%%%%%


\hypertarget{dma}{}
\chapter{通用 DMA 控制器 (GDMA)}
\label{mod:dma}


\section{概述}

通用直接存储访问 (General Direct Memory Access, GDMA) 用于在外设与存储器之间以及存储器与存储器之间提供高速数据传输。软件可以在无需 CPU 干预的情况下通过 GDMA 快速搬移数据，从而降低了 CPU 的工作负载，提高了效率。

GDMA 有 6 个独立的通道，其中包括 3 个发送通道和 3 个接收通道。这 6 个通道被支持 GDMA 功能的外设所共享，用户可以将通道分配给任何支持 GDMA 功能的外设。GDMA 支持的外设包括 PARLIO、ADC、UHCI、AES、SHA、I2S、通用 SPI (GP-SPI)。\chipname{} 的 GDMA 可以访问片内存储器和片外存储器。

GDMA 控制器支持通道间可配置优先级及权重仲裁以管理外设不同的带宽需求。


\begin{figure}[H]
\fbox{\includegraphics[width=0.5\textwidth]{\modulefiles/figures/gdma_ahb_modules_c5.png}}
\centering
\caption{共用 GDMA 通道的模块}
\label{Figure:Module_with_DMA_AHB}
\end{figure}



\section{特性}

GDMA 具有以下特性：

\begin{itemize}
    \item AHB 总线架构
    \item 数据传输以字节为单位，传输数据量可软件编程
    \item 支持任意地址和大小 (size) 访问
    \item 对齐要求：
        \begin{itemize}
          \item 描述符存储地址：字对齐
          \item 数据地址和长度：无要求
        \end{itemize}
    \item 支持链表
    \item 访问存储器时，支持 INCR 突发传输
    \item 三个传输通道和三个接收通道
    \item 可配置支持 GDMA 的外设连接任一通道
    \item 支持通道间优先级和权重仲裁配置
    \item 支持存储搬运功能
\end{itemize}

\section{架构}

\chipname{} 中所有需要进行高速数据传输的模块都具有 GDMA 功能。GDMA 控制器与 CPU 的数据总线使用相同的地址空间访问存储器。图 \ref{fig:gdma-controller} 为 GDMA 控制器基本架构图。

\begin{figure}[H]
\includegraphics[width=1\textwidth]{\modulefiles/figures/gdma_engine.pdf}
\centering
\caption{GDMA 控制器的架构}
\label{fig:gdma-controller}
\end{figure}


GDMA 控制器有 6 个独立的通道，其中包括 3 个发送通道和 3 个接收通道。每个通道可选择与不同的外设相连，从而实现通道资源由外设共享。

GDMA 通过 AHB\_BUS 将数据存入存储器，或者将数据从存储器取出。在通过总线传输数据之前，GDMA 采用可配置的仲裁机制对每个通道的读写请求进行仲裁。存储器的具体使用范围详见章节 \ref{mod:sysmem} \textit{\nameref{mod:sysmem}}。

软件可以通过挂载链表的方式来使用 GDMA 控制器。GDMA 链表存储在存储器中。这些链表包括 outlink\textcolor{red}{\it{n}} 与 inlink\textcolor{red}{\it{n}}，本文以 \textcolor{red}{\it{n}} 来表示通道号，\textcolor{red}{\it{n}} 为 0 \verb+~+ 2。GDMA 从存储器中取得链表，然后根据 outlink\textcolor{red}{\it{n}} 中的内容将相应存储器中的数据发送出去，也可根据 inlink\textcolor{red}{\it{n}} 中的内容将接收的数据存入指定存储器地址空间。


\section{功能描述}

\subsection{链表}

\begin{figure}[H]
\includegraphics[width=0.9\textwidth]{\modulefiles/figures/link_strut}
\centering
\caption{链表结构图}
\label{fig:gdma-link-descriptor}
\end{figure}

图 \ref{fig:gdma-link-descriptor} 所示为链表的结构图。发送链表与接收链表结构相同。每个链表由一个或者若干个描述符构成，一个描述符由 3 个字组成。链表应存放在存储器中，供 GDMA 控制器使用。描述符每一字段的意义如下：

\begin{itemize}
    \item owner (DW0) [31]：表示当前描述符对应的 buffer 允许的操作者。\\
    0：允许的操作者为 CPU。\\
    1：允许的操作者为 GDMA 控制器。\\
    在 GDMA 使用完该描述符对应的 buffer 后，对于 接收描述符，硬件默认会自动将该位清零；对于 发送描述符，需要将 \hyperref[fielddesc:AHBDMAOUTAUTOWRBACKCHN]{AHB\_DMA\_OUT\_AUTO\_WRBACK\_CH\regindex{n}} 置 1，硬件才会自动将该位清零。 软件也可通过置位 \hyperref[fielddesc:AHBDMAOUTLOOPTESTCHN]{AHB\_DMA\_OUT\_LOOP\_TEST\_CH\regindex{n}} 或 \hyperref[fielddesc:AHBDMAINLOOPTESTCHN]{AHB\_DMA\_IN\_LOOP\_TEST\_CH\regindex{n}} 来关闭硬件自动清零的功能。软件在挂载链表时需要将该位置 1。

    {\bfseries 注意： } 本文以 AHB\_DMA\_OUT 开头的寄存器对应 TX 通道寄存器，以 AHB\_DMA\_IN 开头的寄存器对应 RX 通道寄存器。
    \item suc\_eof (DW0) [30]：表示一个描述符对应的数据成功传输后是否触发 AHB\_DMA\_IN\_SUC\_EOF\_CH\regindex{n}\_INT 或 AHB\_DMA\_OUT\_EOF\_CH\regindex{n}\_INT 中断。\\
    0：当前描述符成功传输后不会触发中断；\\
    1：当前描述符成功传输后触发中断。\\
    对于接收描述符，需要软件将该位写 0，硬件会在接收到包含 EOF 标志的数据后将该位置 1。\\对于发送描述符，需要软件按需要将描述符中的该位置 1。
    如果软件在某个描述符中将该位配置为 1，则 GDMA 在处理完该描述符时，会在发送给外设的数据中加入 EOF 标志，告知外设该段数据是一个阶段性结束。
    \item Reserved (DW0) [29]：保留。此位为无关项。
    \item err\_eof (DW0) [28]：表示接收结束错误标志。\\
    0：接收数据没有错误。\\
    1：接收数据有错误。\\
    该位只用于 UHCI 或 PARLIO 利用 GDMA 接收数据。对于接收描述符，硬件在收完描述符对应的数据段并检测到接收数据错误会将该位置 1。
    \item Reserved (DW0) [27:24]：保留。
    \item length (DW0) [23:12]：表示当前描述符对应的 buffer 中的有效字节数。对于发送描述符，该段由软件填写，表示从 buffer 中读取数据时需要读取的字节数；对于接收描述符，该段由硬件使用完该 buffer 后或者接收到最后一个数据时自动填写，表示 buffer 中存储的有效字节数。
    \item size (DW0) [11:0]：表示当前描述符对应的 buffer 容量的字节数。
    \item buffer address pointer (DW1)：buffer 的地址。
    \item next descriptor address (DW2)：下一个描述符的地址。如果当前描述符为链表中最后一个描述符时，该值可以为 0。该地址必须指向存储器的地址空间。
    \begin{tiplisting}
        需要注意，GDMA 的链表描述符本身存储地址要求 4 字节对齐。
    \end{tiplisting}
\end{itemize}


用 GDMA 接收数据时，如果接收数据的长度小于当前描述符指定的 buffer size，那么下一个描述符对应的接收数据不会占用该 buffer 的剩余空间。如果接收数据的长度大于当前描述符指定的 buffer size，那么 GDMA 完成当前描述符后继续获取下一个描述符来完成数据传输。


\subsection{外设到存储器及存储器到外设的数据传输}
GDMA 支持存储器到外设及外设到存储器的数据传输，分别对应 TX 及 RX 功能。TX 通道通过 outlink\textcolor{red}{\it{n}} 实现将指定存储器中的数据搬运到外设的发送端；RX 通道通过 inlink\textcolor{red}{\it{n}} 实现将外设接收到的数据搬运到存储器中。

每个 RX/TX 通道均可以被配置连接到任意一个支持 GDMA 功能的外设，表 \ref{tab:gdma-ahb-peri-sel} 所示为通过寄存器配置通道与其对应外设的关系。其中 ``Dummy-\regindex{n}" 是 GDMA 在实现存储器到存储器数据搬运功能时可选配的寄存器值。当其中一个通道已经与某一个外设连接时，其他通道将不能配置为与该外设连接。

\begin{longtable}[c]{|l|p{3cm}|}
\caption{GDMA 配置寄存器与外设选择关系表}
\label{tab:gdma-ahb-peri-sel}\\
\hline
    \rowcolor{lightgray}
    \textbf{\hyperref[fielddesc:AHBDMAPERIINSELCHN]{AHB\_DMA\_PERI\_IN\_SEL\_CH\regindex{n}}} & \\
    \rowcolor{lightgray}
    \textbf{\hyperref[fielddesc:AHBDMAPERIOUTSELCHN]{AHB\_DMA\_PERI\_OUT\_SEL\_CH\regindex{n}}} & \multirow{-2}{*}{\textbf{外设}}\\\hline

    0   & Dummy-0 \\ \hline
    1   & GP-SPI \\ \hline
    2   & UHCI  \\ \hline
    3   & I2S \\ \hline
    4 \verb+~+ 5 & Dummy-4 \verb+~+ 5 \\ \hline
    6   & AES   \\ \hline
    7   & SHA   \\ \hline
    8   & ADC  \\ \hline
    9   & PARLIO  \\ \hline
    10 \verb+~+  15   & Dummy-10  \verb+~+  15  \\ \hline
    16 \verb+~+  63 & 无效值 \\\hline
\end{longtable}


\subsection{存储器到存储器数据传输}
GDMA 支持存储器到存储器的数据传输。置位 \hyperref[fielddesc:AHBDMAMEMTRANSENCHN]{AHB\_DMA\_MEM\_TRANS\_EN\_CH\regindex{n}}，TX 通道 \textcolor{red}{\it{n}} 的输出将与 RX 通道 \textcolor{red}{\it{n}} 的输入相连，从而使能存储器到存储器的数据传输功能。需要注意的是，一个 TX 通道只与其编号对应的 RX 通道相连而实现存储器到存储器的数据传输，并且需要将 \hyperref[fielddesc:AHBDMAPERIINSELCHN]{AHB\_DMA\_PERI\_IN\_SEL\_CH\regindex{n}} 和 \hyperref[fielddesc:AHBDMAPERIOUTSELCHN]{AHB\_DMA\_PERI\_OUT\_SEL\_CH\regindex{n}} 配置成相同且对应 Dummy 的值。

\subsection{启动 GDMA}

软件通过挂载链表的方式来使用 GDMA。对于接收数据，软件挂载好接收链表并准备好接收数据，配置 \hyperref[fielddesc:AHBDMAINLINKADDRCHN]{AHB\_DMA\_INLINK\_ADDR\_CH\regindex{n}} 字段指向第一个接收链表描述符，置位 \hyperref[fielddesc:AHBDMAINLINKSTARTCHN]{AHB\_DMA\_INLINK\_START\_CH\regindex{n}} 位启动 GDMA。对于发送数据，软件挂载好发送链表并准备好发送数据，配置 \hyperref[fielddesc:AHBDMAOUTLINKADDRCHN]{AHB\_DMA\_OUTLINK\_ADDR\_CH\regindex{n}} 字段指向第一个发送链表描述符，置位 \hyperref[fielddesc:AHBDMAOUTLINKSTARTCHN]{AHB\_DMA\_OUTLINK\_START\_CH\regindex{n}} 位启动 GDMA。\hyperref[fielddesc:AHBDMAINLINKSTARTCHN]{AHB\_DMA\_INLINK\_START\_CH\regindex{n}} 与 \hyperref[fielddesc:AHBDMAOUTLINKSTARTCHN]{AHB\_DMA\_OUTLINK\_START\_CH\regindex{n}} 位由硬件自动清零。

有时您可能想要在 DMA 数据传输已经开始后追加更多描述符。要挂载更多描述符，原本看似只需清空已挂载链表最后一个描述符的 EOF 位，并将该描述符的 next descriptor address (DW2) 字段配置为新链表第一个描述符的地址。但如果 DMA 数据传输已经或即将结束，这个方法便无法奏效。GDMA 控制器有专门的逻辑来确保数据传输继续或重启：如果数据传输仍在进行，GDMA 控制器会确保顾及到新追加的描述符；如果数据传输已经结束，GDMA 控制器会重启数据传输，传输新追加的描述符。这个逻辑由 Restart 功能实现。

软件使用 Restart 功能时，需要重写已挂载链表的最后一个描述符，使其第三个字中的内容（即 DW2）指向新链表的首地址；然后置位 \hyperref[fielddesc:AHBDMAINLINKRESTARTCHN]{AHB\_DMA\_INLINK\_RESTART\_CH\regindex{n}} 或者 \hyperref[fielddesc:AHBDMAOUTLINKRESTARTCHN]{AHB\_DMA\_OUTLINK\_RESTART\_CH\regindex{n}}（这两个位由硬件自动清零），如图 \ref{fig:gdma-link-restart} 所示，硬件会在读取已挂载链表的最后一个描述符时，获取新挂载链表的地址，从而继续处理新挂载的链表。

\begin{figure}[H]
\includegraphics[width=1\textwidth]{\modulefiles/figures/link_restart}
\centering
\caption{链表关系图}
\label{fig:gdma-link-restart}
\end{figure}


\subsection{读链表}

软件在配置并启动 GDMA 后，GDMA 会从存储器中读取链表。GDMA 会检查读入的链表描述符是否正确。只有当链表描述符通过检查时，GDMA 对应的通道 才会开始搬运数据。当链表描述符没有通过检查，硬件将触发描述符错误中断（AHB\_DMA\_IN\_DSCR\_ERR\_CH\regindex{n}\_INT 或者 AHB\_OUT\_DSCR\_ERR\_CH\regindex{n}\_INT），同时该通道将会处于阻塞状态，停止工作。

描述符检查项包括：
\begin{itemize}
    \item \hyperref[fielddesc:AHBDMAINCHECKOWNERCHN]{AHB\_DMA\_IN\_CHECK\_OWNER\_CH\regindex{n}} 或者 \hyperref[fielddesc:AHBDMAOUTCHECKOWNERCHN]{AHB\_DMA\_OUT\_CHECK\_OWNER\_CH\regindex{n}} 置 1 时，检查描述符的 owner 位。如果该位为 0，表示当前操作者为 CPU，则不通过检查。\hyperref[fielddesc:AHBDMAINCHECKOWNERCHN]{AHB\_DMA\_IN\_CHECK\_OWNER\_CH\regindex{n}} 或者 \hyperref[fielddesc:AHBDMAOUTCHECKOWNERCHN]{AHB\_DMA\_OUT\_CHECK\_OWNER\_CH\regindex{n}} 置 0 时，不检查描述符的 owner 位。\\

    \item 检查描述符本身的存储地址是否在配置的存储器地址空间内。如果 GDMA 的描述符地址在 \hyperref[fielddesc:AHBDMAACCESSINTRMEMSTARTADDR]{AHB\_DMA\_ACCESS\_INTR\_MEM\_START\_ADDR} \verb+~+ \hyperref[fielddesc:AHBDMAACCESSINTRMEMENDADDR]{AHB\_DMA\_ACCESS\_INTR\_MEM\_END\_ADDR} 值的范围内，则通过检查。请参见章节 \ref{sec:dma-access-inter-ram}）。\\
\end{itemize}

软件在检查到通道描述符错误中断后，需要准备新描述符，并置位 \hyperref[fielddesc:AHBDMAOUTLINKSTARTCHN]{AHB\_DMA\_OUTLINK\_START\_CH\regindex{n}} 或者  \hyperref[fielddesc:AHBDMAINLINKSTARTCHN]{AHB\_DMA\_INLINK\_START\_CH\regindex{n}} 位启动 GDMA。


\subsection{数据传输结束标志}

{\bfseries 注意：} 本章中提到发送链表描述符的 EOF 为 描述符中的第一个字中的 bit 30 指示的 suc\_eof，接收链表描述符的 EOF 可以为 bit 30 指示的 suc\_eof 和 bit 28 指示的 err\_eof。

GDMA 通过 EOF 来指示对应描述符所需传输的数据段传输结束。

发送数据时，产生 EOF 和 TOTAL EOF 两种中断：
\begin{itemize}
    \item AHB\_DMA\_OUT\_EOF\_CH\regindex{n}\_INT：链表中的任意一个描述符有配置 suc\_eof 标志位，则当该描述符对应的数据传输完成后，GDMA 都会产生 EOF 中断。置位 \hyperref[fielddesc:AHBDMAOUTEOFCHNINTENA]{AHB\_DMA\_OUT\_EOF\_CH\regindex{n}\_INT\_ENA} 位使能该中断。
    \item AHB\_DMA\_OUT\_TOTAL\_EOF\_CH\regindex{n}\_INT：如果该链表中的最后一个描述符配置了 suc\_eof 标志位，则当该链表的最后一个描述符对应的数据传输完成后，GDMA 会产生 TOTAL EOF 中断。置位 \hyperref[fielddesc:AHBDMAOUTTOTALEOFCHNINTENA]{AHB\_DMA\_OUT\_TOTAL\_EOF\_CH\regindex{n}\_INT\_ENA} 位使能该中断。
\end{itemize}


接收数据时，可以产生以下中断：
\begin{itemize}
    \item AHB\_DMA\_IN\_SUC\_EOF\_CH\regindex{n}\_INT：表示带有 EOF 标志的数据段接收完成。置位 \hyperref[fielddesc:AHBDMAINSUCEOFCHNINTENA]{AHB\_DMA\_IN\_SUC\_EOF\_CH\regindex{n}\_INT\_ENA} 位使能该中断。
    \item（仅 UHCI 和 PARLIO）AHB\_DMA\_IN\_ERR\_CH\regindex{n}\_EOF\_INT：表示对应描述符所需传输的数据段接收完成但接收数据有错误。置位 \hyperref[fielddesc:AHBDMAINERREOFCHNINTENA]{AHB\_DMA\_IN\_ERR\_EOF\_CH\regindex{n}\_INT\_ENA} 使能该中断。需要注意的是，只有当通道连接的外设为 UHCI 或 PARLIO 时，才支持该中断。
\end{itemize}


软件在检测到 \hyperref[int:GDMAOUTTOTALEOFINT]{AHB\_DMA\_OUT\_TOTAL\_EOF\_CH\regindex{n}\_INT} 或 \hyperref[int:GDMAINSUCEOFINT]{AHB\_DMA\_IN\_SUC\_EOF\_CH\regindex{n}\_INT} 中断时，可以通过获取 \hyperref[fielddesc:AHBDMAOUTEOFDESADDRCHN]{AHB\_DMA\_OUT\_EOF\_DES\_ADDR\_CH\regindex{n}} 或 \hyperref[fielddesc:AHBDMAINSUCEOFDESADDRCHN]{AHB\_DMA\_IN\_SUC\_EOF\_DES\_ADDR\_CH\regindex{n}} 字段中的的值，获取当前已经完成数据传输的描述符的存储地址。这样，软件可以知道哪些描述符已经被使用并根据需要回收描述符。RX 方向还支持记录 \hyperref[int:GDMAINERREOFINT]{AHB\_DMA\_IN\_ERR\_EOF\_CH\regindex{n}\_INT} 中断触发时描述符的存储地址，可通过 \hyperref[fielddesc:AHBDMAINERREOFDESADDRCHN]{AHB\_DMA\_IN\_ERR\_EOF\_DES\_ADDR\_CH\regindex{n}} 字段获取。


\subsection{访问存储器}\label{sec:dma-access-inter-ram}
GDMA 任意 RX/TX 通道均可以访问存储器，其可访问的地址空间为寄存器 \hyperref[fielddesc:AHBDMAACCESSINTRMEMSTARTADDR]{AHB\_DMA\_ACCESS\_INTR\_MEM\_START\_ADDR} \verb+~+ \hyperref[fielddesc:AHBDMAACCESSINTRMEMENDADDR]{AHB\_DMA\_ACCESS\_INTR\_MEM\_END\_ADDR} 的配置值区间。

为加速数据传输速率，支持突发传输模式。GDMA 默认不使能突发传输模式，用户可以通过配置 \hyperref[fielddesc:AHBDMAINDATABURSTMODESELCHN]{AHB\_DMA\_IN\_DATA\_BURST\_MODE\_SEL\_CH\regindex{n}} 和 \hyperref[fielddesc:AHBDMAOUTDATABURSTMODESELCHN]{AHB\_DMA\_OUT\_DATA\_BURST\_MODE\_SEL\_CH\regindex{n}} 为 0、1、2 将 AHB 总线突发传输长度分别配置为 SINGLE、INCR4 或 INCR8。

GDMA 访问配置的存储器地址空间时，链表描述符参数无对齐要求。

\subsection{仲裁}

为了确保及时响应通用 SPI 等高速低延迟外设的请求，GDMA 引入通道优先级仲裁和通道权重仲裁两种仲裁机制。

\begin{itemize}
\item 优先级仲裁方式：GDMA 支持 6（0 \verb+~+ 5）个等级的优先级，数字越大优先级越高。若多个通道在同一时刻存在数据传输需求，GDMA 会按照通道优先级从高到低顺序响应每个通道的数据传输请求；其数值越大，对应的优先级越高，请求响应越及时
\item 权重仲裁方式：
    \begin{itemize}
        \item 每个通道的的权重（即令牌数）可以配置为 0 \verb+~+ 15 之间的任意值。GDMA 将时间按照一定的总线 (AHB) 时钟周期长度划分成片，在每个时间片内，各个通道的数据传输请求次数应该将按照通道所配置的令牌数进行分配，在任意一个时间片内，某通道每请求一次数据传输消耗一次令牌，若某个通道在某个时间片内消耗尽令牌数，则该通道请求将不再被响应，直到仲裁时间达到时间片规定长度或者所有通道的令牌数均已经消耗尽，此时将刷新令牌数，进入下一个时间片周期内。
        \item 若通道令牌数配置为非零值，即使无数据传输请求，也会在每个时间片内等待所有通道的令牌数耗尽，不会提前退出当前仲裁时间片。通道无数据传输请求且令牌值非 0 时，会导致一些带宽浪费，基于此，GDMA 提供额外的权重仲裁优化开关。权重仲裁优化打开时，如果通道无数据传输需求，则该通道令牌数会被认为是 0，不再参与仲裁，当所有通道均无数据传输请求时，仲裁器提前退出当前时间片，以此提高仲裁响应时间，提高传输速率。
    \end{itemize}
\end{itemize}

\begin{tiplisting}
\vspace{-2em}
\begin{itemize}
    \item 若同时配置优先级仲裁和通道权重仲裁，且优先级不同，权重相同，则在各个通道传输相同的数据量，且数据量较少的前提下，优先级高的通道优先传输完成。
    \item 若优先级不同且权重也不同，各个通道传输数据量相同且较大，则权重较大的通道是能优先分配到大带宽的，会优先完成传输。
\end{itemize}
\end{tiplisting}


\subsection{总线错误响应信息记录}\label{sec:dma-bus-error-resp}

在 GDMA 搬运数据期间，如果总线上从机侧回复了错误响应 (Error Response)，数据传输不会中断，GDMA 依旧会按照已有的链表搬运数据搬运。这样可能会导致完成数据搬运后，实际得到的数据与预期不符。为了便于定位问题，GDMA 传输数据过程中，如果发生了错误相应，可以通过以下内部寄存器存储相关传输信息：

\begin{itemize}
    \item \hyperref[fielddesc:AHBDMAAHBINFRESPERRADDR]{AHB\_DMA\_AHBINF\_RESP\_ERR\_ADDR}：记录本次总线传输对应的 32 位地址。
    \item \hyperref[fielddesc:AHBDMAAHBINFRESPERRCHID]{AHB\_DMA\_AHBINF\_RESP\_ERR\_CH\_ID}：记录本次总线传输使用的 GDMA 通道。
    \item \hyperref[fielddesc:AHBDMAAHBINFRESPERRID]{AHB\_DMA\_AHBINF\_RESP\_ERR\_ID}：记录本次总线上外设与存储器传输或存储器与存储器传输的传输 ID。
    \item \hyperref[fielddesc:AHBDMAAHBINFRESPERRWR]{AHB\_DMA\_AHBINF\_RESP\_ERR\_WR}：记录本次总线传输对应的传输类型。
\end{itemize}

GDMA 作为主机发起 AHB 总线传输且收到错误响应时，还会产生错误相应中断来提醒 CPU。中断产生后，GDMA 仍不会中断数据传输。此外，GDMA 不会锁存第一笔发生错误响应的 AHB 总线传输相关寄存器信息。如果连续发生多笔错误响应传输，上述错误响应记录寄存器的内容会实时更新，最终寄存器中保留的信息是查询寄存器之前最后一笔发生错误响应的 AHB 总线传输。

\subsection{自动时钟门控}\label{sec:dma-auto-clk-gating}

GDMA 在工作时，某些场景只需要用到部分通道，此时可关闭未使用通道的时钟来降低功耗。同理，内部仲裁和总线接口模块仅在有数据传输时才需要工作，在空闲时也可以通过关闭时钟来降低功耗。目前 GDMA 内部有以下模块：
\begin{itemize}
    \item 总线接口处理模块
    \item 仲裁模块
    \item RX 数据处理模块
    \item RX 描述符处理模块
    \item TX 数据处理模块
    \item TX 描述符处理模块
    \item 寄存器配置同步模块
\end{itemize}

这些模块可以在空闲时自动关闭模块时钟，来更大限度的减小功耗，也支持通过寄存器配置强制保持时钟开启。例如，置位 \hyperref[fielddesc:AHBDMAAHBINFCLKEN]{AHB\_DMA\_AHBINF\_CLK\_EN} 可使总线接口处理模块时钟保持常开，其余模块时钟常开配置与其类似。


\section{事件任务矩阵功能} \label{sec:dma-etm-func}

在 \chipname{} 中，GDMA 支持 ETM 功能，即可以通过任意外设的 ETM 事件触发 GDMA 的 ETM 任务，或者通过 GDMA 的 ETM 事件触发任意外设的 ETM 任务。关于 ETM 更多详细信息，请参考章节 \ref{mod:evntaskmatrix} \nameref{mod:evntaskmatrix}。这里仅介绍与 GDMA 相关的 ETM 任务和 ETM 事件。


GDMA 可接收的 ETM 任务有：

\begin{itemize}
\item GDMA\_TASK\_IN\_START\_CH\regindex{n}：触发时开启 GDMA 的接收通道 \regindex{n} 进行数据传输。
\item GDMA\_TASK\_OUT\_START\_CH\regindex{n}：触发时开启 GDMA 的发送通道 \regindex{n} 进行数据传输。
\end{itemize}

\begin{tiplisting}
以上两个 ETM 任务与 CPU 配置 \hyperref[fielddesc:AHBDMAINLINKSTARTCHN]{AHB\_DMA\_INLINK\_START\_CH\regindex{n}}、\hyperref[fielddesc:AHBDMAOUTLINKSTARTCHN]{AHB\_DMA\_OUTLINK\_START\_CH\regindex{n}} 实现的功能相同。 当 \hyperref[fielddesc:AHBDMAINETMENCHN]{AHB\_DMA\_IN\_ETM\_EN\_CH\regindex{n}} 或 \hyperref[fielddesc:AHBDMAOUTETMENCHN]{AHB\_DMA\_OUT\_ETM\_EN\_CH\regindex{n}} 为 1 时，只有 ETM 可以配置对应方向及对应通道 GDMA 的启动，当 \hyperref[fielddesc:AHBDMAINETMENCHN]{AHB\_DMA\_IN\_ETM\_EN\_CH\regindex{n}} 或 \hyperref[fielddesc:AHBDMAOUTETMENCHN]{AHB\_DMA\_OUT\_ETM\_EN\_CH\regindex{n}} 为 0 时，仅可通过 CPU 使能对应的 GDMA。
\end{tiplisting}

GDMA 可产生的 ETM 事件有：

\begin{itemize}
\item GDMA\_EVT\_IN\_DONE\_CH\regindex{n}：表示通道 \regindex{n} 接收方向链表描述符指向的数据接收完成。
\item GDMA\_EVT\_IN\_SUC\_EOF\_CH\regindex{n}：通道 \regindex{n} 接收方向 EOF 位为 1 的链表描述符指向的数据接收完成指示。
\item GDMA\_EVT\_IN\_FIFO\_EMPTY\_CH\regindex{n}：通道 \regindex{n} 接收方向 FIFO 非空到空状态改变指示。
\item GDMA\_EVT\_IN\_FIFO\_FULL\_CH\regindex{n}：通道 \regindex{n} 接收方向 FIFO 非满到满状态改变指示。
\item GDMA\_EVT\_OUT\_DONE\_CH\regindex{n}：通道 \regindex{n} 表示通道 \regindex{n} 发送方向链表描述符指向的数据发送完成。
\item GDMA\_EVT\_OUT\_EOF\_CH\regindex{n}：通道 \regindex{n} 发送方向 EOF 位为 1 的链表描述符指向的数据发送完成指示。
\item GDMA\_EVT\_OUT\_TOTAL\_EOF\_CH\regindex{n}：通道 \regindex{n} 发送方向最后一个并且 EOF 位为 1 的链表描述指向的数据发送符完成指示。
\item GDMA\_EVT\_OUT\_FIFO\_EMPTY\_CH\regindex{n}：通道 \regindex{n} 发送方向 FIFO 非空到空状态改变指示。
\item GDMA\_EVT\_OUT\_FIFO\_FULL\_CH\regindex{n}：通道 \regindex{n} 发送方向 FIFO 非满到满状态改变指示。
\end{itemize}

在具体应用中，GDMA 的 ETM 事件可以用来触发 GDMA 的 ETM 任务，例如，GDMA\_EVT\_OUT\_TOTAL\_EOF\_CH0 事件可以触发 GDMA\_TASK\_IN\_START\_CH1 任务，从而触发新一轮的 GDMA 操作。


\section{中断}

\chipname{} 的 GDMA 可以生成以下中断信号并发送给\hyperref[mod:intmtrx]{中断矩阵}。

    \begin{itemize}
        \item GDMA\_IN\_CH0\_INTR
        \item GDMA\_IN\_CH1\_INTR
        \item GDMA\_IN\_CH2\_INTR
        \item GDMA\_OUT\_CH0\_INTR
        \item GDMA\_OUT\_CH1\_INTR
        \item GDMA\_OUT\_CH2\_INTR
    \end{itemize}

这些中断信号由 GDMA 的内部中断源生成。

以下中断源生成 GDMA\_IN\_CH\regindex{n}\_INTR (\regindex{n} 为 0 \verb+~+ 2) 中断信号：

\begin{itemize}
    \item \label{int:GDMAINDONEINT}AHB\_DMA\_IN\_DONE\_CH\regindex{n}\_INT：对于接收通道\textcolor{red}{\it{n}}，当一个接收链表描述符对应的数据接收完成时触发此中断。
    \item \label{int:GDMAINSUCEOFINT}AHB\_DMA\_IN\_SUC\_EOF\_CH\regindex{n}\_INT：对于接收通道\regindex{n}，当一个接收链表描述符对应的数据接收完成，并且描述符的 suc\_eof 为 1 时触发此中断。
    \item \label{int:GDMAINERREOFINT}AHB\_DMA\_IN\_ERR\_EOF\_CH\regindex{n}\_INT：对于接收通道\textcolor{red}{\it{n}}，当接收的描述符对应数据段中有错误发生时触发此中断。(该中断只用于外设选择 UHCI (UART0/UART1) 或 PARLIO 时）。
    \item \label{int:GDMAINDSCREMPTYINT}AHB\_DMA\_IN\_DSCR\_EMPTY\_CH\regindex{n}\_INT：对于接收通道\textcolor{red}{\it{n}}，当接收链表描述符指向的 buffer size 小于待接收数据长度时触发此中断。
    \item \label{int:GDMAINDSCRERRINT}AHB\_DMA\_IN\_DSCR\_ERR\_CH\regindex{n}\_INT：对于接收通道\textcolor{red}{\it{n}}，当接收链表描述符未通过任一描述符检查时触发此中断。
    \item \label{int:GDMAINFIFOOVFINT} AHB\_DMA\_INFIFO\_OVF\_CH\regindex{n}\_INT：当 GDMA 接收数据方向缓存溢出时触发此中断。
    \item \label{int:GDMAINFIFOUDFINT}AHB\_DMA\_INFIFO\_UDF\_CH\regindex{n}\_INT：当 GDMA 接收数据方向缓存无数据时触发此中断。
    \item \label{int:GDMAINAHBINFRESPERRINT}AHB\_DMA\_IN\_AHBINF\_RESP\_ERR\_CH\regindex{n}\_INT：GDMA 接收通道\textcolor{red}{\it{n}} 进行数据传输时，AHB 总线上如果有错误相应，触发此中断。
\end{itemize}

以下中断源生成 GDMA\_OUT\_CH\regindex{n}\_INTR (\regindex{n} 为 0 \verb+~+ 2) 中断信号：

\begin{itemize}
    \item \label{int:GDMAOUTDONEINT}AHB\_DMA\_OUT\_DONE\_CH\regindex{n}\_INT：对于发送通道\textcolor{red}{\it{n}}，当一个发送链表描述符对应的数据发送完成时触发此中断。
    \item \label{int:GDMAOUTEOFINT}AHB\_DMA\_OUT\_EOF\_CH\regindex{n}\_INT：对于发送通道\textcolor{red}{\it{n}}，当发送描述符的 EOF 位为 1，且 \hyperref[fielddesc:AHBDMAOUTEOFMODECHN]{AHB\_DMA\_OUT\_EOF\_MODE\_CH\regindex{n}} 为 0 时，该描述符对应的最后一个数据由总线传入 GDMA TX 通道时，该中断触发；当发送描述符的 EOF 位为 1，且 \hyperref[fielddesc:AHBDMAOUTEOFMODECHN]{AHB\_DMA\_OUT\_EOF\_MODE\_CH\regindex{n}} 为 1 时，该描述符对应的最后一个数据从 GDMA TX 通道发出时，该中断触发。
    \item \label{int:GDMAOUTDSCRERRINT}AHB\_DMA\_OUT\_DSCR\_ERR\_CH\regindex{n}\_INT：对于发送通道\textcolor{red}{\it{n}}，当发送链表描述符未通过任一描述符检查时会触发此中断。
    \item \label{int:GDMAOUTTOTALEOFINT}AHB\_DMA\_OUT\_TOTAL\_EOF\_CH\regindex{n}\_INT：对于发送通道\textcolor{red}{\it{n}}，当一个链表（可包含多个链表描述符）对应的所有数据都已发送完成时触发此中断。
    \item \label{int:GDMAOUTFIFOOVFINT} AHB\_DMA\_OUTFIFO\_OVF\_CH\regindex{n}\_INT：当 GDMA 发送数据方向缓存溢出时触发此中断。
    \item \label{int:GDMAOUTFIFOUDFINT}AHB\_DMA\_OUTFIFO\_UDF\_CH\regindex{n}\_INT：当 GDMA 发送数据方向缓存无数据时触发此中断。
    \item \label{int:GDMAOUTAHBINFRESPERRINT}AHB\_DMA\_OUT\_AHBINF\_RESP\_ERR\_CH\regindex{n}\_INT：GDMA 发送通道\textcolor{red}{\it{n}} 进行数据传输时，AHB 总线上如果有错误相应，触发此中断。
\end{itemize}


\section{编程流程}

可通过配置 \hyperref[fielddesc:PCRGDMACLKEN]{PCR\_GDMA\_CLK\_EN} 寄存器管理 GDMA 时钟门控，默认时钟门控打开。可通过配置 \hyperref[fielddesc:PCRGDMARSTEN]{PCR\_GDMA\_RST\_EN} 对 GDMA 模块全局复位。

\subsection{GDMA TX 通道配置流程}
利用 GDMA 发送数据时，GDMA TX 通道的软件配置流程如下：
\begin{enumerate}
    \item 对寄存器 \hyperref[fielddesc:AHBDMAOUTRSTCHN]{AHB\_DMA\_OUT\_RST\_CH\regindex{n}} 置 1 然后置 0，复位 GDMA TX 通道状态机和 FIFO 指针；
    \item 挂载好发送链表，配置寄存器 \hyperref[fielddesc:AHBDMAOUTLINKADDRCHN]{AHB\_DMA\_OUTLINK\_ADDR\_CH\regindex{n}} 指向第一个发送链表描述符；
    \item 配置 \hyperref[fielddesc:AHBDMAPERIOUTSELCHN]{AHB\_DMA\_PERI\_OUT\_SEL\_CH\regindex{n}} 为对应的外设号，见表 \ref{tab:gdma-ahb-peri-sel}；
    \item 置位 \hyperref[fielddesc:AHBDMAOUTLINKSTARTCHN]{AHB\_DMA\_OUTLINK\_START\_CH\regindex{n}} 启动 GDMA TX 通道发送数据；
    \item 配置对应的外设，并启动该外设，具体配置请参考对应的外设章节；
    \item 等待 AHB\_DMA\_OUT\_TOTAL\_EOF\_CH\regindex{n}\_INT 中断，即数据传输完成。
\end{enumerate}

\subsection{GDMA RX 通道配置流程}
利用 GDMA 接收数据时，GDMA RX 通道的软件配置流程如下：
\begin{enumerate}
    \item 对寄存器 \hyperref[fielddesc:AHBDMAINRSTCHN]{AHB\_DMA\_IN\_RST\_CH\regindex{n}} 置 1 然后置 0，复位 GDMA RX 通道状态机和 FIFO 指针；
    \item 挂载好接收链表，配置寄存器 \hyperref[fielddesc:AHBDMAINLINKADDRCHN]{AHB\_DMA\_INLINK\_ADDR\_CH\regindex{n}} 指向第一个接收链表描述符；
    \item 配置 \hyperref[fielddesc:AHBDMAPERIINSELCHN]{AHB\_DMA\_PERI\_IN\_SEL\_CH\regindex{n}} 为对应的外设号，见表 \ref{tab:gdma-ahb-peri-sel}；
    \item 置位 \hyperref[fielddesc:AHBDMAINLINKSTARTCHN]{AHB\_DMA\_INLINK\_START\_CH\regindex{n}} 启动 GDMA RX 通道发送数据；
    \item 配置对应的外设，并启动该外设，具体配置请参考对应的外设章节；
\end{enumerate}

\subsection{GDMA 存储器到存储器配置流程}
利用 GDMA 从存储器到存储器搬运数据时配置流程如下：
\begin{enumerate}
    \item 对寄存器 \hyperref[fielddesc:AHBDMAOUTRSTCHN]{AHB\_DMA\_OUT\_RST\_CH\regindex{n}} 置 1 然后置 0，复位 GDMA TX 通道状态机和 FIFO 指针；
    \item 对寄存器 \hyperref[fielddesc:AHBDMAINRSTCHN]{AHB\_DMA\_IN\_RST\_CH\regindex{n}} 置 1 然后置 0，复位 GDMA RX 通道状态机和 FIFO 指针；
    \item 挂载好发送链表，配置寄存器 \hyperref[fielddesc:AHBDMAOUTLINKADDRCHN]{AHB\_DMA\_OUTLINK\_ADDR\_CH\regindex{n}} 指向第一个发送链表描述符；
    \item 挂载好接收链表，配置寄存器 \hyperref[fielddesc:AHBDMAINLINKADDRCHN]{AHB\_DMA\_INLINK\_ADDR\_CH\regindex{n}} 指向第一个接收链表描述符；
    \item 置位 \hyperref[fielddesc:AHBDMAMEMTRANSENCHN]{AHB\_DMA\_MEM\_TRANS\_EN\_CH\regindex{n}} 使能 存储器到存储器 传输功能；
    \item 置位 \hyperref[fielddesc:AHBDMAOUTLINKSTARTCHN]{AHB\_DMA\_OUTLINK\_START\_CH\regindex{n}} 启动 GDMA TX 通道发送数据；
    \item 置位 \hyperref[fielddesc:AHBDMAINLINKSTARTCHN]{AHB\_DMA\_INLINK\_START\_CH\regindex{n}} 启动 GDMA RX 通道发送数据；
    \item 如果某个发送链表描述符的 suc\_eof 位配置为 1，则该描述符对应的数据段传输结束后会触发 AHB\_DMA\_IN\_SUC\_EOF\_CH\regindex{n}\_INT 中断。
\end{enumerate}

\subsection{通道优先级和权重配置流程}

优先级仲裁配置方法如下：
\begin{enumerate}
\item  配置 \hyperref[fielddesc:AHBDMATXPRICHN]{AHB\_DMA\_TX\_PRI\_CH\regindex{n}} 和 \hyperref[fielddesc:AHBDMARXPRICHN]{AHB\_DMA\_RX\_PRI\_CH\regindex{n}} 为 TX 和 RX 通道的优先级。
\end{enumerate}

权重仲裁配置方法如下:
\begin{enumerate}
\item 配置 \hyperref[fielddesc:AHBDMAARBTIMEOUT]{AHB\_DMA\_ARB\_TIMEOUT} 为权重仲裁的时间片长度；
\item 配置 \hyperref[fielddesc:AHBDMATXCHARBWEIGHCHN]{AHB\_DMA\_TX\_CH\_ARB\_WEIGH\_CH\regindex{n}} 和 \hyperref[fielddesc:AHBDMARXCHARBWEIGHCHN]{AHB\_DMA\_RX\_CH\_ARB\_WEIGH\_CH\regindex{n}} 为各个通道的令牌值；
\item 清除 \hyperref[fielddesc:AHBDMATXARBWEIGHOPTDIRCHN]{AHB\_DMA\_TX\_ARB\_WEIGH\_OPT\_DIR\_CH\regindex{n}}，使能 TX 和 RX 方向的权重仲裁优化；
\item 置位 \hyperref[fielddesc:AHBDMAWEIGHTEN]{AHB\_DMA\_WEIGHT\_EN}，使能权重仲裁。
\end{enumerate}


%       for Sections Register Summary and Registers below <<<

%%%%%%%%%%%%%%%%%%%%%%%%
%%% Register Summary %%%
%%%%%%%%%%%%%%%%%%%%%%%%

% >>> Update the sample text below and insert
%     the info on register summary into the subfile <<<

\newpage
\begin{landscape}
\hypertarget{dma-reg-summ}{}
\section{寄存器列表}
\label{sec:dma-reg-summ}

本小节的所有地址均为相对于 GDMA 基地址的地址偏移量（相对地址），具体基地址请见章节 \ref{mod:sysmem} \textit{\nameref{mod:sysmem}} 中的表 \ref{tab:sysmem-base-address}。

请查看章节 \hyperref[glossary-access-types]{\textit{寄存器的访问类型}}，了解“访问”列缩写的含义。

\subfile{\modulefiles/04-DMA-reg-summ__CN}

\end{landscape}

%%%%%%%%%%%%%%%%%
%%% Registers %%%
%%%%%%%%%%%%%%%%%

% >>> Update the sample text below and insert
%      the info on registers into the subfile <<<

\newpage
\section{寄存器}

本小节的所有地址均为相对于 GDMA 基地址的地址偏移量（相对地址），具体基地址请见章节 \ref{mod:sysmem} \textit{\nameref{mod:sysmem}} 中的表 \ref{tab:sysmem-base-address}。

关于保留 (reserved) 域的处理，请查看章节 \hyperref[programming-reserved-field]{如何配置寄存器的保留域}。

\subfile{\modulefiles/04-DMA-reg__CN}

\end{document}
